% 第九卷 4.2.1 节：ChatML 式对话模板的词元序列与损失掩码。
% 三行依次为系统、用户、助手消息；每个方框是一个词元或一段教学用的词元组（中文正文按片段合并画出）。
% 方框下方的数字为损失掩码 m_t；黄色框为 m_t=1 的位置，灰色描边框为特殊标记，白框为普通文本。
\begin{tikzpicture}[
  x=1mm,y=1mm,
  every node/.style={font=\figurefont\footnotesize,text=FigureInk},
  tok/.style={rectangle,rounded corners=1.2pt,draw=FigureGrid,line width=0.6pt,fill=white,
    minimum height=6mm,inner sep=1mm,anchor=west},
  spec/.style={tok,draw=FigureStroke,fill=FigurePaper,font=\figurefont\scriptsize\ttfamily},
  sup/.style={tok,draw=FigureYellow,line width=0.8pt,fill=FigureYellowFill},
  supspec/.style={spec,draw=FigureYellow,line width=0.8pt,fill=FigureYellowFill},
  m/.style={font=\figurefont\scriptsize,text=FigureMuted,anchor=north},
  role/.style={font=\figurefont\scriptsize\bfseries,text=FigureMuted,anchor=east}
]
  % 系统消息。
  \node[role] at (-2,0) {系统};
  \node[spec] (s1) at (0,0) {<|im\_start|>};
  \node[spec] (s2) at ([xshift=1mm]s1.east) {system};
  \node[spec] (s3) at ([xshift=1mm]s2.east) {\textbackslash n};
  \node[tok]  (s4) at ([xshift=1mm]s3.east) {只依据给出的通知回答。};
  \node[spec] (s5) at ([xshift=1mm]s4.east) {<|im\_end|>};
  \foreach \n in {s1,s2,s3,s4,s5} {\node[m] at ([yshift=-0.5mm]\n.south) {0};}

  % 用户消息。
  \node[role] at (-2,-14) {用户};
  \node[spec] (u1) at (0,-14) {<|im\_start|>};
  \node[spec] (u2) at ([xshift=1mm]u1.east) {user};
  \node[spec] (u3) at ([xshift=1mm]u2.east) {\textbackslash n};
  \node[tok]  (u4) at ([xshift=1mm]u3.east) {通知：……外部人员须于16日18时前报名。};
  \node[tok]  (u5) at ([xshift=1mm]u4.east) {外部人员怎么参加？};
  \node[spec] (u6) at ([xshift=1mm]u5.east) {<|im\_end|>};
  \foreach \n in {u1,u2,u3,u4,u5,u6} {\node[m] at ([yshift=-0.5mm]\n.south) {0};}

  % 助手消息：内容与结束标记计分，角色头不计分。
  \node[role] at (-2,-28) {助手};
  \node[spec] (a1) at (0,-28) {<|im\_start|>};
  \node[spec] (a2) at ([xshift=1mm]a1.east) {assistant};
  \node[spec] (a3) at ([xshift=1mm]a2.east) {\textbackslash n};
  \node[sup]  (a4) at ([xshift=1mm]a3.east) {须在};
  \node[sup]  (a5) at ([xshift=1mm]a4.east) {16日};
  \node[sup]  (a6) at ([xshift=1mm]a5.east) {18时前};
  \node[sup]  (a7) at ([xshift=1mm]a6.east) {报名。};
  \node[supspec] (a8) at ([xshift=1mm]a7.east) {<|im\_end|>};
  \foreach \n in {a1,a2,a3} {\node[m] at ([yshift=-0.5mm]\n.south) {0};}
  \foreach \n in {a4,a5,a6,a7,a8} {\node[m,text=FigureInk] at ([yshift=-0.5mm]\n.south) {1};}

\end{tikzpicture}
